\documentclass[11pt, a4paper]{article}
\usepackage[utf8]{inputenc}
\usepackage{amsmath, amssymb}
\usepackage{graphicx}
\usepackage[margin=1in]{geometry}
\usepackage{tcolorbox}
\usepackage{fancyhdr}

\definecolor{UNIBBlue}{RGB}{0, 51, 102}

\pagestyle{fancy}
\fancyhf{}
\rhead{\textbf{Optimisasi untuk Teknik Elektro}}
\lhead{Problem Set - Minggu 3}
\cfoot{\thepage}

\begin{document}

\begin{center}
    \Large\color{UNIBBlue}\textbf{PROBLEM SET}\\
    \large\textbf{Minggu 3: Metode Simpleks dan Dualitas}
\end{center}

\vspace{0.5cm}

\textbf{Instruksi:} Kerjakan soal-soal di bawah ini dengan menyertakan langkah-langkah analitis.

\section*{Soal 1: Pemahaman Konsep (C1 - C2)}
1. (C1) Apakah yang dimaksud dengan variabel slack dan variabel surplus dalam optimisasi linier? \\
2. (C2) Jelaskan hubungan antara Solusi Optimal Primal dan Solusi Optimal Dual (Strong Duality Theorem).

\section*{Soal 2: Formulasi Bentuk Standar (C3)}
Diberikan masalah program linier untuk desain \textit{Microgrid} berikut: \\
Minimumkan $C = 10P_g + 2P_w$ \\
s.t. \\
$P_g + P_w \ge 100$ (Pemenuhan Beban Minimum) \\
$P_g \le 80$ (Kapasitas Grid) \\
$2P_g + 5P_w \le 250$ (Batas Emisi dan Kebisingan Ekuivalen) \\
$P_g, P_w \ge 0$

Ubahlah model di atas ke dalam \textbf{Bentuk Standar} dengan menambahkan variabel slack dan surplus yang sesuai!

\section*{Soal 3: Transformasi Primal-Dual (C4)}
Berdasarkan model awal (Primal) pada Soal 2 di atas, formulasikan model \textbf{Dual}-nya secara lengkap (Tujuan, Variabel, dan Kendala)!

\section*{Soal 4: Evaluasi Solusi dan Shadow Price (C5 - C6)}
Sebuah pabrik produksi komponen elektronik memodelkan maksimasi profitnya. Hasil analisis optimisasi menggunakan \textit{software} JuMP menunjukkan bahwa \textit{shadow price} untuk:
\begin{itemize}
    \item Kendala Mesin Solder (Kapasitas 500 jam/minggu) = Rp 150.000 / jam
    \item Kendala Mesin Testing (Kapasitas 300 jam/minggu) = Rp 0 / jam
    \item Kendala Suplai Bahan Baku (Kapasitas 1000 unit/minggu) = Rp 25.000 / unit
\end{itemize}

Pabrik memiliki dana investasi terbatas untuk \textit{upgrade} peralatan.
(C5) Berikan argumentasi mesin mana (Solder atau Testing) yang menguntungkan untuk di-upgrade kapasitasnya, atau apakah lebih baik mencari supplier bahan baku tambahan? 
(C6) Berikan saran kebijakan strategis kepada manajer operasional terkait mesin yang nilai \textit{shadow price}-nya Rp 0.

\end{document}
